% ===========================================================================
\section{路径编码与连分数}\label{sec:cf}

由推论 \ref{cor:sbn}，每个正有理数对应一条由 $L$、$R$ 组成的路径。把路径中
\emph{连续同向}的段合并，路径就由一串正整数（段长）唯一确定；本节说明这串整数
恰是该有理数的连分数系数。这是``在 Stern--Brocot 树上查找''与``辗转相除''
是同一件事的两种表述。

\subsection{整段移动}

\begin{lemma}[整段移动的端点公式]\label{lem:block-move}
设当前端点为 $a/b<c/d$，$bc-ad=1$，目标分数 $p/q$ 严格介于其间。若连续向右
移动 $t$ 步，则右端点不动，左端点移动到
\[
  \frac{a+tc}{b+td};
\]
连续向左移动 $t$ 步，则左端点不动，右端点移动到
\[
  \frac{ta+c}{tb+d}.
\]
\end{lemma}

\begin{proof}
对 $t$ 归纳，每次向右移动都把左端点替换为它与右端点的中位分数：
$(a,b)\mapsto(a+c,b+d)$，右端点不变，故 $t$ 次后为 $(a+tc,\,b+td)$。向左同理。
\end{proof}

由此，从根到 $p/q$ 的路径可唯一写成
\[
  R^{t_0}L^{t_1}R^{t_2}\cdots,
  \qquad t_0\ge 0,\ t_1,\ldots,t_{n-1}\ge 1,\ t_n\ge 1
\]
（若第一步向左，则记 $t_0=0$；末段次数必为正）。段长序列 $(t_0,t_1,\ldots,t_n)$
就是路径的\keyword{游程编码}。

\subsection{与连分数的对应}

\begin{theorem}[路径的连分数编码]\label{thm:cf-encoding}
设正有理数 $p/q$ 在 Stern--Brocot 树上的路径游程编码为
$(t_0,t_1,\ldots,t_n)$，则
\[
  \frac pq=[t_0;t_1,\ldots,t_n,1],
\]
其中右边是末尾为 $1$ 的有限连分数。反过来，任给一个末尾为 $1$ 的连分数，其系数
（不计最后的 $1$）依奇偶位给出路径：偶数下标项是向右段的长度，奇数下标项是
向左段的长度。
\end{theorem}

\begin{proof}
按引理 \ref{lem:block-move} 记录每一整段移动后的端点位置。不妨设首段向右
（否则 $t_0=0$，不影响论证）。第 $k$ 段移动后的端点记为 $p_k/q_k$：偶数段向右，
记录的是左端点；奇数段向左，记录的是右端点。再补两个虚拟端点
\[
  \frac{p_{-2}}{q_{-2}}=\frac01,\qquad \frac{p_{-1}}{q_{-1}}=\frac10.
\]
由引理 \ref{lem:block-move}，第 $k$ 段移动 $t_k$ 步把端点从 $p_{k-2}/q_{k-2}$
一侧更新为
\[
  \frac{p_k}{q_k}=\frac{t_k\,p_{k-1}+p_{k-2}}{t_k\,q_{k-1}+q_{k-2}}.
\]
这正是连分数渐近分数的递推关系，配合初始条件即得
\[
  \frac{p_k}{q_k}=[t_0;t_1,\ldots,t_k].
\]
搜索停止时，目标分数是最后两个端点的中位分数：
\[
  \frac pq=\frac{p_n+p_{n-1}}{q_n+q_{n-1}}=[t_0;t_1,\ldots,t_n,1],
\]
末位系数 $1$ 对应``再沿最后方向多走一步取中项''。
\end{proof}

有理数的连分数表示由辗转相除给出，故由 $p/q$ 求路径的计算量与辗转相除相同；
反过来，路径（游程编码）也就是 $p/q$ 的连分数。Stern--Brocot 数系与连分数
表示因此是同一编码的两种读法。

\begin{corollary}[父节点与子节点的连分数表达]\label{cor:cf-parent}
设节点 $x=[t_0;t_1,\ldots,t_n,1]$。则其父节点是``沿最后方向少走一步''的节点：
\[
  \text{父节点}=
  \begin{cases}
    [t_0;t_1,\ldots,t_n-1,1],& t_n\ge 2,\\[2pt]
    [t_0;t_1,\ldots,t_{n-1},1],& t_n=1;
  \end{cases}
\]
两个子节点为 $[t_0;t_1,\ldots,t_n+1,1]$ 与 $[t_0;t_1,\ldots,t_n,1,1]$，
左右归属由 $n$ 的奇偶决定。
\end{corollary}

\begin{proof}
父节点是搜索路径删去最后一步的节点：若末段长 $t_n\ge2$，段长减一；若 $t_n=1$，
末段整段消失，倒数第二段的端点成为新端点，对应删去末两位系数。子节点是搜索
再往下走一层：沿末段方向再走一步（$t_n$ 加一），或换向走一段长度为 $1$ 的新段
（追加系数 $1$）。
\end{proof}

\subsection{无理数与有理逼近}

\begin{remark}[无限路径与有理逼近]
把搜索过程施于正无理数 $x$：比较永不停止，$x$ 对应唯一一条无限长的
$\{L,R\}$ 字符串，亦即一个无限连分数。路径上经过的每个分数都是最简分数，
其分母严格递增，且收敛到 $x$——前缀越长的逼近分母越大、越精确。
\emph{但}路径上的逼近误差并非严格递减，路径上的分数也未必是``分母不超过某界
的最佳逼近''：课堂上已指出反例，对 $x=51/100$，搜索路径经过 $1/2$ 与 $2/3$，而
\[
  \Bigl|\frac{51}{100}-\frac12\Bigr|=\frac1{100}
  <\frac{47}{300}=\Bigl|\frac{51}{100}-\frac23\Bigr|,
  \qquad 2<3,
\]
即 $1/2$ 在``分母更小''与``误差更小''两个指标上同时优于后出现的 $2/3$。故
``路径上的每个分数都在 Pareto 前沿上''这一说法不成立。严格的逼近理论
（最佳逼近与渐近分数的刻画）属于连分数的丢番图逼近理论；用 Stern--Brocot 树
求``分母不超过 $N$ 的最佳逼近''时，应当在搜索停止处比较当前两个端点到 $x$ 的
距离，而不能直接断言路径上的某个节点最优。
\end{remark}
